\section{系统总体设计}

\subsection{架构概述}
系统由\textbf{日志节点采集器（Agent）}与\textbf{日志中心（Center）}组成。Agent 分为网关节点实例与微服务节点实例：前者通过 OpenResty Lua 采集 HTTP 流量日志；后者采集应用访问日志。Center 负责接收 Agent 上传、生成关注清单、下发规则、执行二次过滤并写入 Loki。图~\ref{fig:arch} 给出总体架构。

\begin{figure}[htbp]
  \centering
  \includegraphics[width=0.95\textwidth]{fig1_system_overview.pdf}
  \caption{分布式定向日志采集系统总体架构}
  \label{fig:arch}
\end{figure}

\subsection{工作流程}
系统按以下步骤运行：
\begin{enumerate}
  \item 网关在 \texttt{log\_by\_lua} 阶段采集 URL、状态码、响应时延等字段，写入共享内存缓冲区；
  \item 网关 Agent 周期性拉取流量日志并 gRPC 上传至 Center，Center 缓存于 Redis；
  \item Center 对流量日志执行高价值过滤与 URL 聚类，生成 Top-$K$ 关注清单并 gRPC 下发；
  \item 微服务 Agent 按清单匹配本地应用日志，写入固定缓存块，异步压缩上传；
  \item Center 二次过滤、去重后批量推送至 Loki，供 Grafana 查询。
\end{enumerate}

\subsection{部署模式}
\textbf{Sidecar 模式}：Agent 与业务容器共享网络命名空间（\texttt{network\_mode: service:*}），适用于 Kubernetes 或 Docker Compose。\textbf{混合开发模式}：基础设施容器化，Center/Agent 本地调试。原型部署于 WSL2，网关映射端口 8088（规避 Windows IIS 占用 80 端口）。
